\documentclass[10pt,a4paper]{article}
\usepackage[utf8]{inputenc}
\usepackage[T1]{fontenc}
\usepackage{amsmath,amssymb}
\usepackage{booktabs}
\usepackage[round,authoryear]{natbib}
\usepackage{hyperref}
\usepackage[margin=18mm]{geometry}
\usepackage{xcolor}

\hypersetup{colorlinks=true,linkcolor=blue!60!black,citecolor=green!50!black,urlcolor=blue!70!black}
\setlength{\parskip}{0.4em}
\setlength{\parindent}{0pt}
\sloppy

\title{\vspace{-1.5em}AutoPos --- 3D Anchor Self-Calibration from Inter-Anchor Ranging Alone\\[3pt]
       \large Concept Brief (Paper~A: System)}
\author{Zekai Xiao}
\date{\vspace{-1em}June 2026}

\begin{document}
\maketitle
\vspace{-1.0em}

\textit{Part of a deliberately split two-paper program. This is the \textbf{system} paper
(Paper~A): the working, ground-truth-validated calibration system and the vertical limit it does
not cross. A companion \textbf{mechanism} paper (Paper~B) explains \emph{why} that vertical limit
exists. The two are non-overlapping by design.}

\textbf{Why this matters.} Anchor-based UWB positioning needs accurately known anchor coordinates, and
the systems that reach high accuracy obtain them from external infrastructure --- surveyed anchor
positions or angle-of-arrival base stations \citep{zihajehzadeh_2017,kok_2015} --- exactly the infrastructure missing in ambulatory,
out-of-lab settings. Anchor \emph{self}-calibration breaks that circular dependency. Yet even with anchors in hand the
field stays well short of the bar it sets itself --- the optical ``gold standard'' the review ties to
clinically-relevant decision-making ($\sim$1~cm): across 37 UWB$+$IMU systems the \emph{best} reported
accuracy is only $\sim$5~cm (most $10$--$80$~cm; only $\sim$19\% within $10$~cm)
\citep{yogesh_2023,yogesh_2024}. The shortfall is sharpest on the vertical: only 5 of the 37 attempt 3D at all (the other 32 stay 2D), and even those five take the vertical from
IMU fusion over surveyed anchors \citep{kok_2015}, several adding a biomechanical-model constraint
\citep{zihajehzadeh_2017}. The industrial, aerial side --- which this clinical review excludes --- does
the same: a flying platform resolves height with an onboard laser rangefinder and barometer
\citep{xiang_2025_uav}. Across both directions the vertical is imported from an external aid,
\emph{never} recovered from UWB geometry alone. That is the gap AutoPos targets: a body-worn, clinical
tag can carry no such aid, so its usable \textbf{vertical must come from the UWB geometry itself}.

\textbf{Why pure-UWB $Z$, and why $4{+}4$.} A body-worn tag does carry one inertial sensor --- but no
on-body IMU gives a \emph{drift-free, absolute} vertical on its own: an IMU's height comes from
double-integrating acceleration, so it drifts
at a rate set by sensor grade --- even a premium unit --- and is rescued only by gait/biomechanical
assumptions the review notes fail for the individual, non-cyclic movement of patients
\citep{yogesh_2023}; a barometer stays coarse and drifts indoors. Recovering $Z$ from the UWB geometry
supplies the missing absolute reference and keeps the vertical in the \emph{same} drift-free error
regime that already fixes the horizontal, with no second sensor to register or de-drift. The catch is purely
geometric: with a single coplanar layer the vertical is \emph{unobservable} in the ideal range-only
model --- its inter-anchor Cram\'{e}r--Rao bound is exactly singular, so any apparent height comes from
nonidealities or priors, not calibrated geometry, and is far too inaccurate to use. The minimal remedy is the intuitive one --- four anchors is
the well-known minimum for a 3D fix, so we take that single, horizontally-complete layer and copy it to
a second height offset by a $\Delta z$, giving a balanced $4{+}4$: the smallest vertical baseline that
makes the vertical usable, not an arbitrary denser array nor a coplanar anchor layer plus a $Z$ sensor on the tag. Forcing
UWB onto its weakest axis is also what exposed what we name the \emph{delay--layout coupling}
(companion paper).

\textbf{The gap.} Existing self-calibration falls into two clusters. One is inter-anchor-only but
\emph{2D/coplanar}, or needs a few surveyed reference anchors
(\citealp{de_preter_2019}; \citealp{piavanini_2022}; \citealp{corbalan_2023};
\citealp{ridolfi_2021}). The other reaches \emph{3D} but only by adding an external sensor or mobile
platform --- LiDAR, SLAM, or visual-inertial odometry (\citealp{nguyen_2025}; \citealp{yuan_2024};
\citealp{delama_2025}) --- effectively replacing manual surveying with a robot-assisted survey. The
cell ``3D $+$ inter-anchor ranging only $+$ joint antenna-delay $+$ no known reference'' is empty.
It is empty because it is the ill-conditioned regime the field side-steps: \citet{de_preter_2019}
call 3D self-calibration ``very ill-conditioned'' and drop to 2D with tape-measured heights;
\citet{corbalan_2023} report poor vertical resolution and measure height independently; and commercial
RTLS deployments that raise a few anchors to recover the vertical \emph{survey} those raised positions
rather than self-calibrate them. Every escape route imports a measured height --- the manual step this
work removes.

\textbf{What AutoPos does.} On commodity hardware (8$\times$ DWM1001C at two elevation levels) it
takes bidirectional SS-TWR inter-anchor ranging through robust fusion (MAD/MVUE) to a joint solve of
3D anchor coordinates \emph{and} per-device antenna delays --- no external sensor, no mobile rig, no
surveyed reference --- and then validates the layout against optical ground truth. The novelty is a
\emph{deployment capability and its honest characterization}, not a new estimator: the numerical
components (MDS/SDP initialization, Tukey-IRLS robust fitting, alternating delay estimation) are
established, and joint anchor$+$delay estimation is itself prior art
(\citealp{de_preter_2019}; \citealp{shah_2022}). (The deployed firmware applies one generic
antenna-delay default to every node with no per-device calibration, so the solver's delays are the
residuals that default leaves --- the realistic out-of-the-box condition.) What is new is (i) closing the empty 3D
inter-anchor-only cell on cheap single-antenna modules; (ii) a \emph{balanced} two-layer $4{+}4$
array that, \emph{within reference-free self-calibration, and with all eight anchors participating in each fix}, restores the well-conditioned vertical geometry that
a single coplanar layer (or a thinned anchor set) lacks --- deployed few-anchor raises avoid that wall only
by surveying the raised positions, not self-calibrating them (a self-calibration Fisher analysis gives
a vertical CRLB exactly singular for a coplanar layer and $\sim$63~mm for $4{+}4$ --- an RMS lower bound, the same order as the ${\approx}62$~mm
deployed median vertical error). Two ground-truth-free deployment diagnostics also fall out of the
program (not headline results): the ranging residual is a misleading quality signal under minimal
anchor counts---the self-consistent-but-not-metric effect quantified in Paper~B---and estimate
dispersion across the dynamic anchor-selection set is a candidate bias indicator a fixed anchor set
would hide.

\textbf{Result, by axis (Std vs absolute kept separate).} Repeatability $\sim$5~cm 3D / $\sim$4~cm
vertical; absolute (vs Vicon) $\sim$7~cm 3D / $\sim$6~cm vertical --- UWB-only and self-calibrated.
The \emph{vertical} is the point: ${\approx}6$~cm median absolute from pure UWB, on an axis 32 of 37 reviewed systems
avoid entirely and that the remaining few reach only by importing an external height reference. The defensible claim is a usable vertical recovered from UWB alone under the
body-worn constraint that rules out external height aids --- not the smallest absolute number (the
best fusion-plus-surveyed system still edges the 3D total).

\textbf{Anchor visibility.} At every position $92.2\%$ of frames ranged to all eight anchors and
$99.9\%$ to at least seven (per-position median $8/8$; link-level success $99.0\%$). A seven-of-eight fix is still non-coplanar and over-determined, so the vertical stays
well-conditioned in $99.9\%$ of frames and the reported medians are effectively whole-array fixes.

\textbf{Honest limits.} The preliminary $\sim$5~cm numbers are \emph{repeatability}, not accuracy ---
only the optical-ground-truth numbers are reported as accuracy. The $4{+}4$ geometry is
\emph{necessary but not sufficient} for the vertical: it completes the geometric observability, but
a systematic vertical bias remains --- a \emph{tag-side delay} aliasing into the vertical, not a
geometry deficit --- which the companion mechanism paper (Paper~B) analyzes as the tag-side face of
its delay--layout coupling result. These numbers also presuppose that every fix draws on all eight anchors of the two-layer
array --- it is the anchor participation, not the ranging protocol (sequential and broadcast are
identical for a static fix), that conditions the vertical: a fix from fewer anchors (e.g.\ a four-anchor
subset) hits a much larger vertical CRLB, and a coplanar configuration returns the vertical to the singular regime by construction; neither reproduces these numbers. All numbers are from a single room with 24 static
tag positions; multi-room transfer is unproven. The headline 3D median (${\approx}73$~mm) also benefits, in part, from the
production layout's expanded scale compensating an uncorrected tag-side delay --- the companion paper's
mechanism; with that delay left at zero the metric-correct layout gives a larger median. The absolute vertical is
also referenced to the optical marker, so a fixed marker-to-antenna-phase-centre offset would shift the
absolute value (Paper~B's phase-center sensitivity check finds the layout ranking robust to marker
offsets beyond 10~mm).

\textbf{Scope.} To our knowledge this is the first demonstration that the
\emph{vertical} axis of a UWB anchor array can be self-calibrated to the centimetre level
(${\approx}62$~mm median absolute vs Vicon, against a Fisher-predicted $\sim$63~mm inter-anchor vertical CRLB for $4{+}4$) from
inter-anchor ranging alone on commodity single-antenna modules --- a reference point that later reference-free
pure-UWB vertical self-calibration in this regime can be measured against. The full enabling combination is stated once (3D self-calibration; inter-anchor ranging only; joint per-device
antenna-delay estimation; commodity DWM1001C modules; balanced $4{+}4$ array with all eight anchors participating in each fix; optical-ground-truth
validation --- no survey, no external $Z$ sensor, no mobile platform), and ships the CRLB script and
ranging data so the floor is independently reproducible.

\begin{thebibliography}{99}
\small
\bibitem[Corbalan(2023)]{corbalan_2023} Corbalan, P.\ et~al., 2023. Self-localization of UWB anchors:
From theory to practice. \emph{IEEE Access}, 11, pp.29711--29725.
\bibitem[De Preter(2019)]{de_preter_2019} De Preter, A.\ et~al., 2019. Range bias modeling and
autocalibration of a UWB positioning system. In: \emph{IPIN}, pp.1--8.
\bibitem[Delama(2025)]{delama_2025} Delama, G.\ et~al., 2025. Real-time initialization of unknown
anchors for UWB-aided navigation. In: \emph{IROS}. arXiv:2506.15518.
\bibitem[Kok(2015)]{kok_2015} Kok, M.\ et~al., 2015. Indoor positioning using ultrawideband and
inertial measurements. \emph{IEEE Trans.\ Veh.\ Technol.}, 64(4), pp.1293--1303.
\bibitem[Nguyen(2025)]{nguyen_2025} Nguyen, T.-D.\ et~al., 2025. Coordinate-consistent localization
via continuous-time calibration and fusion of UWB and SLAM. \emph{arXiv}:2510.05992.
\bibitem[Piavanini(2022)]{piavanini_2022} Piavanini, M.\ et~al., 2022. A self-calibrating
localization solution for sport applications with UWB technology. \emph{Sensors}, 22(23), 9363.
\bibitem[Ridolfi(2021)]{ridolfi_2021} Ridolfi, M.\ et~al., 2021. UWB anchor nodes self-calibration in
NLOS conditions. \emph{Wireless Networks}, 27, pp.3007--3023.
\bibitem[Shah(2022)]{shah_2022} Shah, S.\ et~al., 2022. Node calibration in UWB-based RTLSs using
multiple simultaneous ranging. \emph{Sensors}, 22(3), 864.
\bibitem[Xiang(2025)]{xiang_2025_uav} Xiang, Z., Chen, L., Wu, Q., Yang, J., Dai, X.\ and Xie, X.,
2025. An improved UWB indoor positioning approach for UAVs based on the dual-anchor model.
\emph{Sensors}, 25(4), 1052.
\bibitem[Yogesh(2023)]{yogesh_2023} Yogesh, V.\ et~al., 2023. Integrated UWB/MIMU sensor system for
position estimation: A technical review (37 systems). \emph{Sensors}, 23(16), 7277.
\bibitem[Yogesh(2024)]{yogesh_2024} Yogesh, V.\ et~al., 2024. Novel calibration method for improved
UWB sensor distance measurement for 3D analysis of human movement. \emph{Eng.\ Sci.\ Technol.\ Int.\
J.}, 58, 101844.
\bibitem[Yuan(2024)]{yuan_2024} Yuan, S.\ et~al., 2024. Large-scale UWB anchor calibration and
one-shot localization using Gaussian process. \emph{arXiv}:2412.16880.
\bibitem[Zihajehzadeh(2017)]{zihajehzadeh_2017} Zihajehzadeh, S.\ and Park, E.J., 2017. A novel
biomechanical model-aided IMU/UWB fusion for magnetometer-free lower body motion capture. \emph{IEEE
Trans.\ Syst.\ Man Cybern.\ Syst.}, 47(6), pp.927--938.
\end{thebibliography}

\end{document}
